\section{Quantitative names and finite consumers}

\subsection{Rational interval names}

A real-valued observation is represented by a finite certificate containing
rational endpoints \(\ell\le u\) with an explicit width bound.  At precision
\(p\), validity includes
\[
  u-\ell\le 2^{-p}
\]
and a theorem-specific enclosure statement.  Signed integers and rationals are
coded by natural tuples; positivity of denominators and cross-multiplication
are QF arithmetic conditions.

The selected-certificate theorem gives a fixed PR term
\(c(p,q)\).  Define, for example, the midpoint observation
\[
  x_p(q)=\frac{\ell(c(p,q),q)+u(c(p,q),q)}2.
\]
If every valid interval contains the same target in the relevant constructive
sense, then two independent intervals satisfy
\[
  |x_p(q)-x_r(q)|
  \le \frac{\operatorname{width}(c(p,q))+
                 \operatorname{width}(c(r,q))}{2}.
\]
This is a cross-certificate Cauchy estimate.

\subsection{Operations by demand transformers}

An operation on names is a finite contract.  For example, to compute
\(x+y\) within \(2^{-p}\), request each input within \(2^{-(p+1)}\).  For
multiplication, the demand additionally depends on explicit magnitude bounds.
For inversion, it depends on an apartness bound
\(|x|\ge2^{-g}\).  Composition of operations is handled by backward demand.

\begin{interface}[Quantitative operation]
A name operation provides:
\begin{enumerate}
\item a PR demand transformer \(d(p,q)\);
\item a finite PR transducer on input certificates;
\item an object proof that valid input certificates at demand \(d(p,q)\)
produce a valid output certificate at demand \(p\);
\item cross-certificate stability for the output observable.
\end{enumerate}
\end{interface}

The demand need not be optimal.  It must be explicit and proved sufficient.

\subsection{Finite consumers and metastable windows}

A bounded downstream computation may inspect only finitely many approximations.
It does not always require a global Cauchy modulus.  Suppose a consumer asks
queries \(p_0,\ldots,p_{m-1}\) and its correctness is stable whenever all
answers lie in specified tolerances.  Backward demand computes one finite
window of producer requirements.  A theorem providing that window is enough.

\begin{theorem}[Finite-consumer closure]\label{thm:finite-consumer}
Let a fixed bounded consumer use at most \(m(q)\) observations, with every
adaptive query bounded by a PR demand envelope.  If the producer supplies a
valid, mutually stable family of certificates on that finite demand-closed
window, then the consumer's output specification follows.  No global modulus
outside the window is needed.
\end{theorem}

\begin{proof}
Encode the complete interaction transcript as one certificate: consumer state,
queries, producer answers, and local checker results.  The query envelope gives
a numerical code bound.  The transcript checker verifies each transition and
answer validity.  Existence of the finite stable window produces one valid
transcript.  Apply whole-path selection and the terminal extractor.
Alternatively, when the consumer and producer are deterministic fixed PR
programs, use a packed-state trace and QF induction.  In neither proof is any
claim made about precisions the consumer never requests.
\end{proof}

This theorem is the correct setting for metastability and ``finite use''
arguments.  It must not be misread as a proof of a global Cauchy modulus.

\subsection{When a genuine infinite name is justified}

To claim a full quantitative name, the package should provide a PR modulus and
a theorem such as
\[
  p,r\ge M(k,q)	o
  \dist(x_p(q),x_r(q))\le2^{-k}.
\]
Cross-certificate stability often makes this easy once certificate errors have
known rates.  The selected terms \(c(p,q)\) form a PR sequence, so no countable
choice is needed.  What remains essential is the explicit modulus or rate.

\subsection{Functions and higher-type-looking objects}

A function is represented operationally: a query contains an input code and an
output precision.  Uniform continuity, differentiability, or integrability is
expressed through explicit moduli and finite remainder certificates.  A
functional consuming another presented function is permitted only through a
finite query protocol whose use bound is explicit.

Thus the framework can express rich higher mathematics without postulating
higher-type objects inside the first-order theory.  The price---and the
constructive content---is that every finite use is visible.
